\documentclass[12pt]{article}
\usepackage[margin=0.75in]{geometry}
\usepackage{fontspec}
\setmainfont{Latin Modern Roman}
\usepackage{microtype}
\usepackage{multicol}
\usepackage{fancyhdr}
\usepackage{titlesec}
\usepackage{enumitem}
\usepackage{xcolor}
\usepackage[hidelinks]{hyperref}
\setlength{\columnsep}{0.28in}
\setlength{\parindent}{1.1em}
\setlength{\parskip}{0.32em}
\setlength{\headheight}{15pt}
\titleformat{\section}{\large\bfseries\scshape}{}{0pt}{}
\titleformat{\subsection}{\normalsize\bfseries}{}{0pt}{}
\pagestyle{fancy}
\fancyhf{}
\fancyfoot[C]{\thepage}
\renewcommand{\headrulewidth}{0.4pt}
\newcommand{\focusbox}[1]{\par\vspace{0.4em}\noindent\begingroup\setlength{\fboxsep}{6pt}\colorbox{gray!12}{\begin{minipage}{\dimexpr\linewidth-2\fboxsep\relax}#1\end{minipage}}\endgroup\par\vspace{0.5em}}

\fancyhead[L]{Whately: Testimony}
\fancyhead[R]{Student Reading}
\begin{document}
\begin{center}
{\LARGE\bfseries Testimony, Authority, and Credibility}\par\vspace{0.25em}
{\large\scshape Week 13: Whom Should We Trust?}\par\vspace{0.5em}
{Adapted and modernized from Richard Whately, \textit{Elements of Logic}}\par\vspace{0.6em}
\rule{0.82\textwidth}{0.4pt}
\end{center}

\section*{Central Question}
When should we trust a source?

\section*{Vocabulary Preview}
\begin{multicols}{2}
\begin{itemize}[leftmargin=*, itemsep=0.18em]
\item \textbf{Testimony}: a claim we receive from another person, witness, record, or source.
\item \textbf{Authority}: a source treated as worth hearing because of office, training, experience, or access.
\item \textbf{Credibility}: the degree to which a source deserves trust in this case.
\item \textbf{Expertise}: knowledge or skill in the relevant field.
\item \textbf{Access}: whether the source was in a position to know.
\item \textbf{Bias}: a pressure that may tilt the source toward one answer.
\item \textbf{Corroboration}: independent support from other sources or evidence.
\item \textbf{Relevance}: whether the source's authority applies to this exact question.
\item \textbf{Probability}: the strength of support, short of certainty.
\item \textbf{Credibility map}: claim, source, access, expertise, bias, corroboration, and careful conclusion.
\end{itemize}
\end{multicols}

\section*{Introduction}
Week 12 asked how examples support general claims. Week 13 asks how testimony supports belief. Much of what we know comes through other people: books, witnesses, teachers, records, news reports, experts, documents, and ordinary conversation.

Whately does not ask students to despise authority. He also does not ask them to accept authority blindly. A disciplined thinker asks what kind of source is speaking, whether the source could know, whether the source has reason to distort, and whether other evidence confirms the claim.

\focusbox{\textbf{Source note:} Adapted and modernized from Whately's treatment of testimony, probability, and the limits of what logic can decide by form alone. Classroom examples are original. The goal is a reasoning habit for school research, public claims, historical evidence, and literary judgment.}

\begin{multicols}{2}
\section*{Reading}

\subsection*{1. Why Testimony Matters}
A person cannot directly observe everything worth knowing. You trust testimony whenever you learn about an event before you were born, a place you have not visited, a scientific result you have not personally tested, or a classroom procedure announced by someone else.

This dependence is not weakness. It is part of human learning. The question is not whether to trust testimony, but how to judge it well.

\subsection*{2. Authority Is Not Magic}
Authority can be useful evidence. A trained mechanic usually knows more about an engine than a passerby. A witness standing near an accident may know something distant readers do not. A historian who has studied many records may deserve attention.

But authority is not magic proof. The expert may speak outside the field. The witness may have seen only part of the event. The official may have an incentive to protect an institution. The popular source may be confident and still be wrong.

\subsection*{3. Four Questions for a Source}
Use four questions before accepting testimony:
\begin{enumerate}[leftmargin=*,itemsep=0.12em]
\item \textbf{Access:} Was the source in a position to know?
\item \textbf{Competence:} Does the source understand the matter?
\item \textbf{Honesty and bias:} Is there pressure to exaggerate, hide, flatter, sell, or attack?
\item \textbf{Corroboration:} Do independent sources or evidence support the claim?
\end{enumerate}

\subsection*{4. Relevant Expertise}
Expertise must match the question. A coach may be a strong source about training habits, but not automatically about nutrition science. A popular actor may be persuasive about a product, but popularity is not product expertise. A classmate may be reliable about what happened in yesterday's club meeting if the classmate was present, but not about a private meeting the classmate never attended.

\subsection*{5. Bias Does Not Automatically Destroy Testimony}
Bias is a warning sign, not a magic eraser. A biased source may still tell the truth. A source with no obvious bias may still be mistaken. The point is to adjust the degree of trust and ask for corroboration.

For example, a student defending a late assignment has a reason to make the delay sound unavoidable. That does not prove the student is lying. It does mean the claim may need records, dates, or another witness before the teacher treats it as settled.

\subsection*{6. Corroboration}
Corroboration is stronger when support comes from independent sources. Three students repeating the same rumor from one group chat do not count as three independent witnesses. A repair receipt, a timestamped message, and an unrelated witness are stronger because they do not all depend on the same voice.

\subsection*{7. Probability and Careful Conclusions}
Testimony often gives probability, not certainty. Good reasoning states the conclusion at the right strength:
\begin{itemize}[leftmargin=*]
\item weak source, little support: ``possible, but not settled'';
\item relevant source, some support: ``reasonable to believe for now'';
\item expert source plus independent evidence: ``strongly supported'';
\item contradicted or unsupported source: ``needs more investigation.''
\end{itemize}

\subsection*{8. Literary Cases}
In \textit{Macbeth}, the witches have strange access to future events, but their testimony is ambiguous and morally dangerous. In \textit{Julius Caesar}, Rome is full of public claims about honor, ambition, loyalty, and liberty. A disciplined reader asks who is speaking, what the speaker could know, what the speaker wants, and what evidence confirms or limits the testimony.

\subsection*{9. The Credibility Map}
Use this map:
\begin{enumerate}[leftmargin=*,itemsep=0.12em]
\item State the claim being testified to.
\item Identify the source.
\item Check access and competence.
\item Name possible bias or pressure.
\item Look for corroboration or conflict.
\item State a careful conclusion about trust.
\end{enumerate}

\section*{Reading Focus}
\begin{enumerate}[leftmargin=*]
\item Why is testimony necessary for human learning?
\item Why is authority useful but not magic proof?
\item List the four questions for judging a source.
\item What is relevant expertise?
\item Why does bias not automatically prove a source false?
\item What makes corroboration independent?
\item How should a careful conclusion change when evidence is weak, moderate, or strong?
\end{enumerate}

\section*{Handout Summary}
Write 5--7 sentences explaining how to evaluate testimony without either rejecting all authority or accepting authority blindly. Your summary must include the words \textit{access}, \textit{expertise}, \textit{bias}, and \textit{corroboration}.
\end{multicols}
\end{document}
